% ═══════════════════════════════════════════════════════════════════════
\section{L'anatomie d'un document}
% ═══════════════════════════════════════════════════════════════════════

\begin{frame}[fragile]{Trois parties, toujours dans cet ordre}
  \begin{columns}[c]
    \begin{column}{0.6\textwidth}
\begin{CodeTeX}[title={main.tex}]
\documentclass[12pt]{article}   % 1. la classe

% 2. le préambule : les réglages
\usepackage[T1]{fontenc}
\usepackage[french]{babel}
\usepackage{amsmath, amssymb}
\title{Mon premier document}
\author{Prénom Nom}

\begin{document}               % 3. le corps
\maketitle
Le texte et les formules vont ici.
\end{document}
\end{CodeTeX}
    \end{column}
    \begin{column}{0.37\textwidth}
      \small
      \begin{enumerate}
        \item La classe : le type de document.
        \item Le préambule : extensions et réglages. Rien ne s'y
          affiche.
        \item Le corps, entre \cmd{begin\{document\}} et
          \cmd{end\{document\}} : tout ce qui sera imprimé.
      \end{enumerate}
      \vspace{4pt}
      Tout ce qui suit un \texttt{\%} sur une ligne est un commentaire,
      ignoré à la compilation.
    \end{column}
  \end{columns}
\end{frame}

\begin{frame}{Les classes de documents}
  \small \texttt{\textbackslash documentclass[options]\{classe\}} : la
  première ligne décide du type de document.
  \begin{center}
  \begin{tabular}{lll}
    \toprule
    Classe & Pour quoi faire & Exemple en enseignement \\
    \midrule
    \texttt{article} & document court, sections & fiche d'activité,
      plan de leçon \\
    \texttt{report}, \texttt{book} & document long, chapitres & notes de
      cours, cahier d'exercices \\
    \texttt{exam} & questions, points, solutions & examen, quiz,
      devoir \\
    \texttt{beamer} & diapositives & présentation en classe (comme
      celle-ci) \\
    \bottomrule
  \end{tabular}
  \end{center}
  \vspace{4pt}
  Options courantes : taille du texte (\texttt{10pt}, \texttt{11pt},
  \texttt{12pt}), format du papier (\texttt{letterpaper} est le défaut en
  Amérique du Nord), \texttt{twocolumn}.
\end{frame}

\begin{frame}[fragile]{Un préambule en français, et pas plus}
  \begin{columns}[c]
    \begin{column}{0.56\textwidth}
\begin{CodeTeX}[title={Le préambule de nos exemples}]
\usepackage[T1]{fontenc}
\usepackage[french]{babel}
\usepackage{amsmath, amssymb}
\usepackage{siunitx}
\sisetup{output-decimal-marker = {,}}
\usepackage[margin=2.5cm]{geometry}
\usepackage{graphicx}
\usepackage{tikz}
\end{CodeTeX}
    \end{column}
    \begin{column}{0.41\textwidth}
      \footnotesize
      \begin{itemize}
        \item \texttt{fontenc}, \texttt{babel} : accents, césure et
          typographie françaises (espaces avant « : ; ? ! »).
        \item \texttt{amsmath}, \texttt{amssymb} : outils et symboles
          mathématiques.
        \item \texttt{siunitx} : nombres et unités, virgule décimale.
        \item \texttt{geometry} : les marges.
        \item \texttt{graphicx} : insérer des images.
        \item \texttt{tikz} : dessiner des figures.
      \end{itemize}
    \end{column}
  \end{columns}
  \vspace{4pt}
  \begin{Rem}{Conseil}
    \small Un préambule court est un préambule qui compile. Si l'IA vous
    propose quinze extensions, demandez-lui à quoi sert chacune.
  \end{Rem}
\end{frame}

\begin{frame}[fragile]{Commandes et environnements}
  \small Deux formes à reconnaître, partout :
\begin{CodeTeX}[title={Une commande, avec option et argument}]
\commande[option facultative]{argument obligatoire}
\section{Mise en situation}       \usepackage[french]{babel}
\end{CodeTeX}
\begin{CodeTeX}[title={Un environnement : un début et une fin qui se répondent}]
\begin{nom}
   ...contenu...
\end{nom}
\end{CodeTeX}
  \begin{itemize}
    \item Les accolades \texttt{\{ \}} entourent toujours un argument ; elles
      vont par paires.
    \item Les majuscules comptent : \cmd{Section} n'existe pas.
    \item Chaque \cmd{begin\{x\}} doit avoir son \cmd{end\{x\}}.
  \end{itemize}
\end{frame}

\begin{frame}[fragile]{Structurer le texte}
  \begin{columns}[T]
    \begin{column}{0.55\textwidth}
\begin{CodeTeX}[title={Sections et paragraphes}]
\section{Mise en situation}
Un paragraphe se termine par
une ligne vide.

Un nouveau paragraphe.

\subsection{Matériel}
Calculatrice, règle.
\end{CodeTeX}
    \end{column}
    \begin{column}{0.42\textwidth}
      \small
      \begin{itemize}
        \item La numérotation des sections est automatique.
        \item \cmd{section*\{...\}} : une section non numérotée.
        \item Un seul retour à la ligne ne fait rien ; une ligne vide
          commence un paragraphe.
        \item Les espaces multiples comptent pour un seul.
      \end{itemize}
    \end{column}
  \end{columns}
\end{frame}

\begin{frame}[fragile]{Les listes}
\begin{CodeRendu}[title={Numérotée ou à puces}]
\begin{enumerate}
  \item Lire l'énoncé.
  \item Faire un schéma.
\end{enumerate}
\begin{itemize}
  \item une puce
\end{itemize}
\end{CodeRendu}
  \small Chaque élément commence par \cmd{item}. Dans un document de classe
  \texttt{article}, le rendu est plus sobre que dans ces diapositives.
\end{frame}

\begin{frame}{Les caractères réservés}
  \small Dix caractères ont un sens particulier pour \LaTeX. Pour les
  afficher :
  \begin{center}
  \begin{tabular}{lll}
    \toprule
    Caractère & Son rôle & Pour l'écrire \\
    \midrule
    \texttt{\%} & commentaire & \texttt{\textbackslash\%} \\
    \texttt{\$} & mode mathématique & \texttt{\textbackslash\$} \\
    \texttt{\&} & colonnes d'un tableau & \texttt{\textbackslash\&} \\
    \texttt{\#} & paramètre d'une commande & \texttt{\textbackslash\#} \\
    \texttt{\_} \quad \texttt{\^{}} & indice, exposant (en mode math) &
      \texttt{\textbackslash\_} \quad \texttt{\textbackslash\^{}\{\}} \\
    \texttt{\{} \quad \texttt{\}} & arguments & \texttt{\textbackslash\{}
      \quad \texttt{\textbackslash\}} \\
    \texttt{\textbackslash} & début d'une commande &
      \texttt{\textbackslash textbackslash} \\
    \texttt{\~{}} & espace insécable & \texttt{\textbackslash\~{}\{\}} \\
    \bottomrule
  \end{tabular}
  \end{center}
  \begin{AuTableau}[Piège fréquent]
    \small « Une hausse de 15\% » : sans la barre oblique, \LaTeX\
    ignore la fin de la ligne. Le symbole \texttt{\%} est le plus souvent
    oublié.
  \end{AuTableau}
\end{frame}

\begin{frame}{À vous}
  \begin{Exercice}{Modifier à la main}
    \small Dans \texttt{01-minimal.tex} :
    \begin{enumerate}
      \item Passez de \texttt{12pt} à \texttt{11pt} ; observez le PDF.
      \item Ajoutez une section « Conclusion » avec deux paragraphes.
      \item Écrivez la phrase « Le rabais est de 20 \% sur 45 \$. » de
        façon qu'elle s'affiche correctement.
      \item Transformez la liste numérotée en liste à puces.
    \end{enumerate}
  \end{Exercice}
\end{frame}
